\section{Quadratic characters from a finite cover}
\label{sec:quadratic-beta-characters}

The quadratic construction uses a cover of degree thirty-six of the
three-punctured sphere. Its homology separates the cubic character pair
from the other sixth-root characters. An involution on that pair lifts
through every finite central extension, which is the additional fact
needed to pass from the finite homology calculation to a full peripheral
automorphism.

Use the modular marking
\[
 \Gamma=\widehat{C_2*C_3}=\langle s,r\mid s^2=r^3=1\rangle,
 \qquad \tau=sr,
 \qquad F=\ker(\Gamma\to S_3),
\]
and choose $x=\tau^{-2}$ and $y=sxs$ as free generators of $F$.
Set $z=(xy)^{-1}$. Conjugation by $s$ interchanges $x,y$, and
conjugation by $r$ cyclically permutes $x,y,z$. Put
\[
 U_6=\overline{F'F^6},\qquad
 C=F/U_6\simeq(\mathbb Z/6\mathbb Z)^2,
 \qquad \Gamma_6=\Gamma/U_6.
\]
Here $F'$ denotes the closed commutator subgroup and $F^6$ the closed
subgroup generated by sixth powers. The group $\Gamma_6$ has order
$216$ and abelianization $C_6$: the inclusion $U_6\subset\Gamma'$
follows because $x,y$ have order dividing three in $\Gamma^{\mathrm{ab}}$.

\subsection{The homology of the cover}

For a prime $p\geq5$, write
\[
 M_p=U_6/\overline{U_6'U_6^p}=H_1(U_6,\mathbb F_p).
\]
It is a $37$-dimensional $\mathbb F_p$-representation of $\Gamma_6$.
To see its restriction to $C$, use the cellular chain complex of the
degree-thirty-six cover of the two-loop graph. Over a field containing
the sixth roots of unity it is
\[
 0\longrightarrow H_1(U_6)\longrightarrow k[C]^2
   \xrightarrow{\ (X-1,\,Y-1)\ }k[C]
   \longrightarrow k\longrightarrow0.
\]
Consequently the trivial character occurs twice, and each nontrivial
character of $C$ occurs exactly once. At a character $(a,b)\ne(1,1)$,
the commutator $[x,y]$ gives a nonzero vector in this one-dimensional
space: its two edge coordinates are, up to the chosen commutator sign,
$(1-b,a-1)$.

Suppose now that $p\equiv5\pmod{12}$, and let
$\omega\in\mathbb F_{p^2}$ be a primitive cube root. The two
characters
\[
 \chi=(\omega,\omega),\qquad
 \chi^{-1}=(\omega^2,\omega^2)
\]
give a two-dimensional $\mathbb F_p$-summand $V_p\subset M_p$.
It is stable under $\Gamma_6$: both characters are fixed by the
permutations of $x,y,z$, and they are interchanged by Frobenius.
Let $W_p$ be its complementary sum of character spaces, so
$M_p=V_p\oplus W_p$.

On the $\chi$ line, conjugation by $s$ acts by $-1$, while
conjugation by $r$ acts by $\omega$. Indeed
\[
 s[x,y]s^{-1}=[y,x],\qquad
 r[x,y]r^{-1}=[y,(xy)^{-1}],
\]
and the second commutator has coefficient $(XY)^{-1}$ times that
of $[x,y]$ in the metabelian commutator module. Its value at
$(\omega,\omega)$ is $\omega$. Thus the eigenvalues of $s,r,\tau$
on this line are $-1,\omega,-\omega$, respectively; none is one.
The conjugate eigenvalues occur on the $\chi^{-1}$ line.

\begin{lemma}\label{lem:cover-sign-homology}
Let $\mathcal P$ be a finite set of primes congruent to five modulo
twelve, and let
\[
 Q=\Gamma\Big/
       \bigcap_{p\in\mathcal P}\overline{U_6'U_6^p}.
\]
Then $Q$ fits into a split exact sequence
\[
 1\longrightarrow\prod_{p\in\mathcal P}M_p
   \longrightarrow Q\longrightarrow\Gamma_6\longrightarrow1,
\]
and $Q^{\mathrm{ab}}=C_6$. Choose a complement to its kernel.
For $p\in\mathcal P$, let $\alpha_p$ act as $-1$ on $V_p$, and
as the identity on $W_p$, on all $M_q$ with $q\ne p$, and on this
complement. Then $\alpha_p$ is an involutive automorphism of $Q$.
It acts trivially on $H_2(Q,\mathbb Z)$.
\end{lemma}

\begin{proof}
The simultaneous elementary abelian quotients of $U_6$ have product
image because their orders are pairwise coprime. Their kernel is the
displayed intersection, proving the exact sequence. The kernel order
is coprime to $216$, so the extension splits by the coprime splitting
theorem. The summand $V_p$ is invariant, hence its sign commutes with
the complement action and defines the asserted automorphism.
The source epimorphism gives a surjection $C_6=\Gamma^{\mathrm{ab}}
\to Q^{\mathrm{ab}}$, and the quotient to $\Gamma_6$ gives the
inverse lower bound, proving equality.

For a prime $\ell\notin\mathcal P$, the positive homology of the
kernel with $\mathbb Z_\ell$ coefficients vanishes by averaging,
so $H_2(Q,\mathbb Z)_\ell$ is identified with
$H_2(\Gamma_6,\mathbb Z)_\ell$. For $q\in\mathcal P$, remove
all other primary factors first. Since the quotient by $M_q$ has
order prime to $q$, the homology spectral sequence collapses to
\[
 H_2(Q,\mathbb Z)_q
      \simeq\bigl(H_2(M_q,\mathbb Z)\bigr)_{\Gamma_6}
      \simeq\bigl(\textstyle\bigwedge^2M_q\bigr)_{\Gamma_6}.
\]
Here the Schur multiplier of an abelian group is its exterior square;
the averaging contraction gives the spectral-sequence collapse.
The automorphism $\alpha_p$ is the identity on all terms with
$q\ne p$. For $q=p$, decompose
\[
 \bigwedge^2M_p=
   \bigwedge^2V_p\ \oplus\ (V_p\otimes W_p)\ \oplus\ \bigwedge^2W_p.
\]
It acts by $+1$ on the two exterior-square summands and by $-1$ on
the middle summand. That middle summand has zero $C$-coinvariants:
after adjoining the sixth roots, its character could be trivial
only if $W_p$ contained $\chi$ or $\chi^{-1}$, and both are in
$V_p$. Since $|C|=36$ is invertible modulo $p$, invariants and
coinvariants are exact and give the same vanishing. The middle
summand therefore contributes nothing after taking $\Gamma_6$
coinvariants. This proves the assertion about $H_2$.
\end{proof}

\subsection{Lifting the sign through central extensions}

\begin{lemma}\label{lem:sign-central-extensions}
Let $Q$ and $\alpha_p$ be as in
Lemma~\ref{lem:cover-sign-homology}, and let $q:\Gamma\to Q$ be the
natural quotient. Every finite central extension
\[
 1\longrightarrow Z\longrightarrow D\longrightarrow Q
       \longrightarrow1
\]
admits an involutive automorphism $\widetilde\alpha_p$ which lifts
$\alpha_p$ and fixes $Z$ pointwise.
For each $t\in\{s,r,\tau\}$, choose $b_t\in V_p$ such that
\[
 \alpha_p(q(t))=b_tq(t)b_t^{-1}.
\]
Such an element exists. For every lift $\widetilde t\in D$ of
$q(t)$, there is a lift $\widetilde b_t\in D$ of $b_t$ with
\[
 \widetilde\alpha_p(\widetilde t)
      =\widetilde b_t\widetilde t\widetilde b_t^{-1}.
\]
The same assertions hold with $\tau^{-1}$ as the third peripheral.
\end{lemma}

\begin{proof}
Write $Z=Z_p\times Z_{p'}$. First quotient by $Z_{p'}$.
The universal-coefficient exact sequence for trivial coefficients is
\[
 0\longrightarrow
 \operatorname{Ext}^1_{\mathbb Z}(Q^{\mathrm{ab}},Z_p)
 \longrightarrow H^2(Q,Z_p)
 \longrightarrow\operatorname{Hom}(H_2(Q,\mathbb Z),Z_p)
 \longrightarrow0.
\]
The first term is zero because $Q^{\mathrm{ab}}=C_6$ and $p\geq5$.
Lemma~\ref{lem:cover-sign-homology} therefore shows that
$\alpha_p$ fixes every class in $H^2(Q,Z_p)$. It lifts to the
given central extension with identity action on its kernel.
Such a lift is unique: two lifts differ by a homomorphism
$Q\to Z_p$, and $\operatorname{Hom}(C_6,Z_p)=0$.
Its square is consequently the identity.

For the $Z_{p'}$ extension, put $Q_0=Q/M_p$. Its order is prime to
$p$. Positive cohomology of the normal $p$-group $M_p$ with
coefficients in $Z_{p'}$ vanishes by averaging, so inflation gives
\[
 H^2(Q_0,Z_{p'})\xrightarrow{\ \sim\ }H^2(Q,Z_{p'}).
\]
The extension is therefore a pullback $Q\times_{Q_0}D_0$ of a
central extension of $Q_0$. Since $\alpha_p$ is the identity on
$Q_0$, the map $(a,d)\mapsto(\alpha_p(a),d)$ is an involutive
lift fixing $Z_{p'}$. Combining the two primary extensions over
$Q$ gives the required involution of $D$.

To construct $b_t$, use the chosen semidirect-product coordinates
of $Q$. The difference between $\alpha_p(q(t))$ and $q(t)$ lies
in $V_p$ and is minus twice the $V_p$-coordinate of $q(t)$.
Conjugation by $b\in V_p$ changes that coordinate by $(1-t)b$.
The operators $1-s,1-r,1-\tau$ are invertible on $V_p$ by the
eigenvalues computed above, so each equation has a unique solution.

Choose a lift $\widetilde b_t$ whose image in the pullback
$Q\times_{Q_0}D_0$ is $(b_t,1)$. Comparing the two lifts of
$\alpha_p(q(t))$ gives
\[
 \widetilde\alpha_p(\widetilde t)
   =\widetilde b_t\widetilde t\widetilde b_t^{-1}z
       \qquad(z\in Z).
\]
Its image in the $Z_{p'}$ extension shows that $z\in Z_p$.
Furthermore $\widetilde\alpha_p(\widetilde b_t)\widetilde b_t$
lies in the central kernel, because its image in $Q$ is $-b_t+b_t=0$.
Applying $\widetilde\alpha_p$ to the displayed equation and using
its involutivity and its identity action on $Z$ yields
$\widetilde t=\widetilde t z^2$. Squaring is injective on the
finite $p$-group $Z_p$, so $z=1$. The assertion for the inverse
peripheral follows by inversion.
\end{proof}

\begin{theorem}[Full lifts of the finite signs]
\label{thm:quadratic-sign-lifts}
For every finite set $\mathcal P$ of primes congruent to five modulo
twelve and each $p\in\mathcal P$, there is a full multiplier-one
gentle element whose action on the sixth-root homology characters
modulo each $q\in\mathcal P$ is $-1$ on
$(\omega,\omega),(\omega^2,\omega^2)$ when $q=p$, and is the
identity on all other nontrivial characters.
One may at the same time prescribe identity conjugator values in
any fixed finite quotient of $\Gamma$ of order divisible only by
two and three.
\end{theorem}

\begin{proof}
Use the finite quotient $Q$ above. Consider any finite quotient
$E$ of $\Gamma$ equipped with its quotient to $Q$, and put
$A=\ker(E\to Q)$. Lemma~\ref{lem:finite-peripheral-refinement}
constructs a finite central extension $D^\dagger\to Q$ and a
source epimorphism $q^\dagger$. Apply
Lemma~\ref{lem:sign-central-extensions} to obtain
$\widetilde\alpha_p$ and three peripheral conjugators over the
fixed $b_s,b_r,b_{\tau^{-1}}\in V_p$. Postcomposing $q^\dagger$
with this automorphism gives the new epimorphism and conjugators
required by the finite-refinement lemma. It supplies an epimorphism
to $E$ with the prescribed conjugator images in $Q$.

Now impose the peripheral product relation, generation in all finite
quotients of $\Gamma$, and the three prescribed conjugator values in
$Q$. The compactness and Hopf argument in the proof of
Theorem~\ref{thm:peripheral-lifting} applies verbatim: every finite
list is controlled by one joint quotient $E$ as just considered.
We obtain a full peripheral automorphism $\Phi$ of $\Gamma$
inducing $\alpha_p$ on $Q$.

An inner correction by the third full conjugator makes $\Phi(\tau)=\tau$.
That conjugator has image in $V_p$, so this correction acts trivially
on every $M_q$ and on $\Gamma_6$; it preserves all the homology
actions under discussion. Since the resulting automorphism is the
identity modulo $U_6$, write
\[
 \Phi(s)=sf,\qquad \Phi(r)=f^{-1}r,
       \qquad f\in U_6.
\]
The square relation says that conjugation by $s$ takes $f$ to $f^{-1}$.
Thus its image in $F^{\mathrm{ab}}$ has the form $(a,-a)$, with
$a\in6\widehat{\mathbb Z}$. Composing with conjugation by $x^a$
replaces $f$ by $y^af x^{-a}$, whose abelianization is zero.
This correction still fixes $\tau$. It acts trivially on every
$M_q$ because the deck image of $x^a$ in $C$ is the identity.
We have therefore reached $f\in F'$. The relations
$(sf)^2=(f^{-1}r)^3=1$ are exactly the normalized multiplier-one
two- and three-cycle equations. The restriction to $F$ is
\[
 x\longmapsto x,\qquad y\longmapsto f^{-1}yf,
\]
and gives the required gentle element.

For the final assertion, choose a positive integer $m$ divisible by
six and by the exponent of $T$, with no prime factors other than two
and three. Replace $U_6$ by $U_m=\overline{F'F^m}$ in the finite
construction. The deck group is now $(\mathbb Z/m\mathbb Z)^2$,
the base has order $6m^2$, and its homology still has multiplicity
one at every nontrivial character. Thus the same two cubic character
spaces define $V_p$ and all the preceding sign and homology arguments
apply without change. Take the joint image of this quotient and $T$.
Its normal kernel is still the full product of the $M_q$, because
the common quotient with $T$ kills this kernel of coprime order.
The new base is a group of order divisible only by two and three,
and its action on the homology factors through $\Gamma/U_m$.
Choose the $b_t$ in $V_p$, so their images in $T$ are the identity.
The first inner normalization has trivial image in $T$. The second
uses $a\in m\widehat{\mathbb Z}$, so $x^a$ also has trivial image
in $T$. Both therefore preserve the identity frame values.
\end{proof}

\subsection{The unitary relation without a metabelian lifting assumption}

We record the finite-cover argument which supplies the norm-one
condition. It uses only positive peripheral preservation.

\begin{lemma}\label{lem:cover-unitarity}
Let $\phi$ be a continuous automorphism of $F$ inducing the identity
on $F^{\mathrm{ab}}$ and preserving the three positive peripheral
classes. On a nontrivial sixth-root character $(a,b)$ with
\[
 a\ne1,\qquad b\ne1,\qquad ab\ne1,
\]
let $B_\phi(a,b)$ be the scalar by which it acts on the corresponding
one-dimensional summand of $H_1(U_6,\overline{\mathbb F}_p)$,
where $p\geq5$. Then
\[
 B_\phi(a,b)B_\phi(a^{-1},b^{-1})=1.
\]
\end{lemma}

\begin{proof}
The degree-thirty-six cover of the three-punctured sphere has six
boundary components over each of its three boundaries, since the
three deck elements $X,Y,XY$ have order six. It has eighteen
boundaries. Euler characteristic gives genus ten. Its profinite
fundamental group is $U_6$, and it has a surface presentation
\[
 \prod_{j=1}^{10}[a_j,b_j]\prod_{i=1}^{18}c_i=1,
\]
where the $c_i$ are positively oriented boundary elements.
Let $H=H_1(U_6,\mathbb F_p)$ and let $R$ be the span of the
boundary classes. It has dimension seventeen, their only relation
being $\sum_i[c_i]=0$.

The alternating form on $H/R$ can be seen directly in the class-two,
exponent-$p$ quotient of $U_6$. Write this quotient in logarithmic
coordinates, which are valid because $p$ is odd, and project its
degree-two commutators from $\bigwedge^2H$ to
$\bigwedge^2(H/R)$. The boundary relation gives the nondegenerate
bivector
\[
 \Omega=\sum_{j=1}^{10}[a_j]\wedge[b_j].
\]
An automorphism preserving the positive boundary conjugacy classes
permutes the $c_i$ up to conjugation. Conjugating a boundary element
changes its degree-two coordinate by a wedge involving its boundary
class, which is zero after this projection. The boundary factors'
pairwise multiplication terms are zero there for the same reason.
Applying the automorphism to the displayed surface relation therefore
gives $(\bigwedge^2\phi)\Omega=\Omega$. This is precisely preservation
of the associated nondegenerate alternating form on $H/R$.

The automorphism $\phi$ preserves $U_6$ and its positive boundary
classes, possibly permuting the components over each original
boundary. It commutes with the deck group on homology, because it
induces the identity on $F/U_6$. The boundary representation over
each original boundary is induced from the corresponding cyclic
deck subgroup. Thus a character $(a,b)$ occurs in $R$ only when
$a=1$, $b=1$, or $ab=1$; conversely each nontrivial character with
one of these equalities occurs there. The trivial representation
occurs twice in $R$. This also accounts for its dimension seventeen.

The two one-dimensional spaces at $(a,b)$ and $(a^{-1},b^{-1})$
in the statement consequently survive in $H/R$ and pair perfectly.
Deck invariance of the form pairs a character only with its inverse.
Preservation of the form forces the product of the two scalar actions
to be one, as asserted.
\end{proof}

\subsection{The quadratic characters and their independence}

For $g$ in the gentle group, use the normalized Alexander coefficient
$B_g(X,Y)$ in the convention of \cite[\S2.3, p.~101]{Furusho2003}.
At multiplier one it is the scalar just used on every nontrivial
character space. At an arbitrary multiplier $\lambda_g$, its relation
to the semilinear coefficient on the commutator is
\[
 A_g(X,Y)=S_{\lambda_g}(X)S_{\lambda_g}(Y)B_g(X,Y),
 \qquad S_a(T)=\frac{T^a-1}{T-1}.
\]
The divided power is the continuous geometric sum. Its identity
$S_{ab}(T)=S_a(T)S_b(T^a)$, together with composition of the
actions, gives
\begin{equation}
 B_{gh}(X,Y)=B_g(X,Y)
                  B_h(X^{\lambda_g},Y^{\lambda_g}).
 \label{eq:beta-crossed-law}
\end{equation}
For pentagon elements, Ihara's beta factorization gives
\begin{equation}
 B_g(a,b)B_g(ab,c)=B_g(b,c)B_g(a,bc).
 \label{eq:beta-associativity}
\end{equation}
We use here its full profinite scope, as recalled in
\cite[\S2.3, p.~101]{Furusho2003}, following
\cite[Proposition~1.6.1]{Ihara1999}.

\begin{theorem}\label{thm:quadratic-beta-family}
For every prime $p\equiv5\pmod{12}$ there is a continuous character
\[
 \vartheta_p:\widehat{\mathrm{GT}}_0\longrightarrow\mathbb F_2
\]
which vanishes on $\widehat{\mathrm{GT}}$. Their joint map is
surjective, already on the multiplier-one subgroup:
\[
 (\vartheta_p)_p:
 \widehat{\mathrm{GT}}_{0,1}
      \twoheadrightarrow\prod_{p\equiv5\ (12)}\mathbb F_2.
\]
Consequently the continuous quadratic character space of the gentle
group vanishing on the pentagon group, and also the one vanishing on
$G_{\mathbb Q}$, has dimension $\aleph_0$.
\end{theorem}

\begin{proof}
Fix $p$ and a primitive cube root $\omega\in\mathbb F_{p^2}$.
All the following values mean reduction modulo $p$ of the continuous
sixth-root evaluations. Put
\[
 E_p(g)=
 \frac{B_g(\omega,\omega)B_g(\omega^2,-1)}
      {B_g(\omega,-1)B_g(\omega,-\omega)}
        \in\mathbb F_{p^2}^{\times}.
\]
Every denominator is nonzero because these coefficients describe
invertible semilinear actions on the one-dimensional character spaces.
Equation~\eqref{eq:beta-crossed-law} gives
\[
 E_p(gh)=E_p(g)\,\sigma_{\lambda_g}(E_p(h)),
\]
where $\sigma_\lambda$ is the identity if $\lambda\equiv1\pmod3$
and Frobenius if $\lambda\equiv-1\pmod3$. Indeed every multiplier
is odd, so it fixes $-1$ and either fixes or interchanges the two
cube roots. Equation~\eqref{eq:beta-associativity}, with
$(a,b,c)=(\omega,\omega,-1)$, gives $E_p(g)=1$ for every pentagon
element.

All four character pairs in the definition satisfy the nonboundary
conditions of Lemma~\ref{lem:cover-unitarity}. For multiplier one,
that lemma says $E_p(g)E_p(g)^p=1$, since Frobenius is simultaneous
inversion on these sixth roots. For an arbitrary $g$, choose a Galois
element $h$ with the same multiplier and write $g=uh$, where
$\lambda_u=1$. The crossed law and $E_p(h)=1$ give $E_p(g)=E_p(u)$.
Thus every $E_p(g)$ lies in the norm-one subgroup, of order $p+1$.
Define
\[
 (-1)^{\vartheta_p(g)}=E_p(g)^{(p+1)/2}.
\]
This is a sign. Frobenius acts by inversion on the norm-one subgroup,
which fixes its quadratic sign, so the crossed law makes
$\vartheta_p$ an ordinary continuous homomorphism. It vanishes on
the pentagon group. Replacing $\omega$ by $\omega^2$ inverts $E_p$
and leaves the sign unchanged.

For a finite set $\mathcal P$ of these primes, the witness of
Theorem~\ref{thm:quadratic-sign-lifts} has at its selected prime $p$
the four coefficient values $(-1,1,1,1)$, and at every other prime
in $\mathcal P$ the four values $(1,1,1,1)$. Hence its $E_p$ value
is $-1$ and all the other $E_q$ values are one. Since
$p\equiv5\pmod{12}$, the integer $(p+1)/2$ is odd. The character
values on these witnesses form the identity matrix. Every finite
projection of the joint map is therefore onto; its image is compact
and hence closed, proving surjectivity onto the full product.

There are infinitely many primes congruent to five modulo twelve:
the Chebotarev theorem applied to the class $5$ in
$\operatorname{Gal}(\Q(\zeta_{12})/\Q)$ gives density $1/4$
\cite[Chapter V, Theorem 3.23, p.~164]{MilneCFT}.
The family therefore has countably infinite
rank. Conversely the marked gentle group is a compact metrizable
subspace of $\widehat{\mathbb Z}^{\times}\times\widehat F_2$.
It has countably many clopen sets, and consequently only countably
many continuous maps to $\mathbb F_2$. This proves the upper bound
for both character spaces. The Galois assertion uses
$G_{\mathbb Q}\subset\widehat{\mathrm{GT}}$.
\end{proof}

\begin{corollary}[A simultaneous profinite abelian quotient]
\label{cor:simultaneous-quotient}
There is a continuous surjection, already onto on global multiplier one,
\[
 \left(\GTz/\overline{\langle\!\langle\GT\rangle\!\rangle}\right)^{\mathrm{ab}}
 \longrightarrow
 (\Z/2\Z)^{\mathbb N}\times
 \prod_{\substack{p\geq5\\p\ \mathrm{prime}}}
       \prod_{k,j\geq1}\Z/p^k\Z.
\]
Consequently its continuous character group with values in $\Q/\Z$
contains the direct sum of the dual finite cyclic groups displayed here.
\end{corollary}
\begin{proof}
Combine Theorem~\ref{thm:quadratic-beta-family} with
Corollary~\ref{cor:all-orders}.  Every finite projection is a subdirect
product of finite primary groups of pairwise coprime orders.  Such a
subdirect product is the full product: its order is divisible by the
order of every primary factor, hence by their product.  Compactness gives
surjectivity of the infinite map.  Pullback of the finite-coordinate
characters gives the stated embedding of their direct sum.
\end{proof}
